\section{Conclusion}
\label{sec:conclusion}

When the unit of evaluation is the laundering attempt, the numbers AMLworld produces are
governed first by choices that papers rarely state and only then by the method. The scope rule
is worth 0.35--0.38 ring-F1. The standard split's tail inflates edge metrics and halves the
effective seed budget in the two days that carry ordinary traffic. And 44\,\% of in-scope
rings cannot be matched as connected objects by any method, because the split and the window
sever them.
% source: docs/FINDINGS.md (C9, C1 and dated notes on the window-matched budget and component-level scope)

Within those limits, detector quality decides ring recovery only up to a point. On HI-Small,
three graph neural detectors from edge F1 0.596 to 0.715 give seeds-only ring-F1 0.518 to 0.551,
and nine attempts to spend the budget better or reassemble rings after the detector leave
ring-F1 flat; most remaining losses are rings with too few visible members. Below that range the
detector still matters: on HI-Medium a detector trained only on HI-Small, applied without
retraining, gives $0.396 \pm 0.015$ against 0.090 for an earlier in-domain detector that differs
in architecture and training data, mostly through ranking. TAML, converted to rings, reaches
0.327 and loses mainly through what its graphs omit.
% source: results/notes/final_numbers.md (detector swap, HI-Medium transfer, TAML); results/notes/recovery_attempts.md

The practical consequence is for how group-level AML results are reported. A ring-level number
is interpretable only with its scope rule, split, window, budget allocation and matching
threshold attached, and edge F1 is not evidence about it. We release the protocol, the tables
behind every number in this paper, and the code.
